\section{Exact phase certificates for short-period overlaps}
\label{sec:lcs-periodic}

\subsection{Recognizing a structured case before general matching}

The full-overlap control of Section~\ref{sec:lcs-windows} asks for all
positive $k$ satisfying
\[
 \operatorname{suffix}_k((ab)^N)=\operatorname{prefix}_k((ab)^N).
\]
Its answer is simply the progression
$2,4,\ldots,2N$, but the general periodic-extension machinery still
exhausts its allowance on the larger test sizes. A short known repeating
block gives a more direct route. It can be verified throughout a compressed
grammar by a finite-state computation, without occurrence tables or long
substring equality queries.

The new stage attempts this route only when both feasible overlap operands
have length at least 128. It reads a prefix of exactly 64 letters from the
right operand and proposes its least period $p\le32$, if one exists.
These fixed constants bound the speculative work and the size of phase
masks; they are policy choices, not mathematical thresholds for general
recognition. The proposed block is $B=b_0\cdots b_{p-1}$.
The sample does not prove that the remaining word is periodic. Both entire
operands must pass the grammar certificate below. Otherwise the existing
complete matcher runs unchanged.

This is another instance of adaptive implementation selection: a small
structural test can choose a cheaper exact algorithm. Failure of that test
is not failure of the underlying search, and exhaustion remains inconclusive.

\subsection{A phase mask at every grammar rule}

For a grammar word $W$, define
\[
 M(W)=\{r\in\mathbb Z/p\mathbb Z:
          W_i=b_{(r+i)\bmod p}\text{ for every }0\le i<|W|\}.
\]
This set is represented by a $p$-bit integer. A terminal $a$ has precisely
the phases at which $b_r=a$. At a binary rule $W=UV$,
\[
 M(UV)=M(U)\cap\bigl(M(V)-|U|\bigr).
\]
Equivalently, rotate the right child's bit mask to the right by
$|U|\bmod p$, then intersect it with the left mask. This condition is
necessary and sufficient: the first child starts at phase $r$ and the
second at phase $r+|U|$. Induction over the grammar proves that every
mask describes the entire represented word, including positions far beyond
the sampled prefix. Signed letters are compared exactly, with no hashing.

The certificate succeeds only if $0\in M(Y)$ and $M(X)$ is nonempty.
Select any certified starting phase $r\in M(X)$ and set
$e=(r+|X|)\bmod p$. The right operand is then a prefix of $B^\infty$,
and every suffix of $X$ starts at its certified phase $e-k$.

\subsection{Recovering every overlap length}

The candidate block is cyclically primitive. Indeed, if a nonzero cyclic
shift fixed $B$, it would be a repetition of a block of length
$\gcd(p,r)<p$. The sampled prefix, being a prefix of $B^\infty$, would
then have that smaller period, contradicting the chosen least period.
Consequently a full block matches $B$ only at phase zero.

For $k\ge p$ an overlap therefore exists exactly when
\[
 k\equiv e\pmod p.
\]
The same congruence is sufficient for shorter lengths as well. It produces
one arithmetic progression with first element $e$ if $e>0$, or $p$ if
$e=0$, and step $p$, clipped at $\min(|X|,|Y|)$. Its multiplicity is stored
as an integer even when exponentially large.

Shorter overlaps outside this congruence class are possible because a
proper fragment of a primitive block may agree at several phases.
The implementation checks each $1\le k<p$ outside the main progression
by comparing
\[
 b_{(e-k+i)\bmod p}=b_i\qquad(0\le i<k).
\]
Each successful check contributes a singleton progression. This gives at
most $p$ disjoint progressions, all exact. Their union is the complete
overlap set, rather than a selected useful witness. The existing critical
alignment routine can therefore consume them without changes to its proof
or the ordinary relator certificate format.

For $X=Y=(ab)^N$, $p=2$ and $e=0$, so the output is $(2,2,N)$.
For a block of 32 distinct letters, with $X$ starting at phase seven and
both words containing $N$ full blocks, the answer is $(7,32,N)$. The latter
case checks a nonzero rotation and the maximal supported period.

\subsection{Work, memory and the scope of the bound}

The prefix sample contains at most 64 letters; its explicit local traversal
has bounded size, independent of the represented word lengths. Creating
that prefix slice costs $O(H)$ for grammar height $H$. Period proposal
and short-exception checks take bounded work because $p\le32$.
The mask pass uses one bounded-width integer per reachable rule and one
transition per binary rule. Its work is $O(s)$ after a topological order
is available. The current shared reachability helper sorts node identifiers,
so an implementation-level bound is $O(s\log s+H)$ word operations,
plus polynomial costs for the binary-encoded lengths. No loop counts to
$N$ or enumerates the members of a returned progression.

The arena's shared allowance charges sample traversal, candidate-period
checks, reachability, mask evaluation and short-overlap checks. Existing
node and overlap-cache limits remain in force. At most one mask is retained
per reachable node, and each mask has at most 32 bits. A failed proposal
adds bounded sample work; a plausible sample that fails globally adds one
grammar pass before the fallback. The negative controls measure this
speculation cost.

This stage establishes an efficient exact local algorithm for two words
certified to follow a common short periodic block. It does not bound the
periods arising throughout unknot recognition, guarantee that relator search
reaches such states, or bound the number of moves and intermediate grammar
sizes. The general subexponential and quasi-polynomial targets remain open
for this implementation.

\subsection{Validation and matched measurements}

Literal-oracle tests cover signed blocks, nonprimitive supplied spellings,
random grammar parses, phase rotations, truncated final blocks, and
exceptional short overlaps. Mutations after position 64 and a 33-letter
period verify that a plausible prefix or an unsupported period does not
bypass the complete matcher. A 32-letter shifted block repeated $2^{500}$
times must finish under 20000 work units with whole-word expansion and occurrence
tables forbidden. A separate oracle checks that the returned progressions
feed the exact critical-extension calculation correctly.
The integration run passes 693 tests in 183.274 seconds. The four-test
periodic suite, including the additional extension oracle, passes in
0.786 seconds; together the runs cover all 694 tests currently present.

\begin{center}\small
\begin{tabular}{@{}lrrr@{}}
\toprule
Scope & Before & A/A & Phases \\
\midrule
Small 14, summed medians & 300.269 & 309.512 & 299.232 \\
Gordian, full query & 2113.390 & 2228.724 & 2140.719 \\
\bottomrule
\end{tabular}
\end{center}
Median milliseconds; all full knot queries complete.

\begin{center}\small
\begin{tabular}{@{}llrrr@{}}
\toprule
Family & $h$ & Before & A/A & Phases \\
\midrule
Prefix & 8 & 0.111 & 0.112 & 0.111 \\
Prefix & 32 & 0.277 & 0.277 & 0.283 \\
Prefix & 500 & 3.609 & 3.483 & 3.577 \\
Interior & 8 & 0.164 & 0.166 & 0.163 \\
Interior & 32 & 0.406 & 0.404 & 0.406 \\
Interior & 500 & 5.254 & 5.227 & 5.184 \\
Equal pairs & 8 & 2.230 & 2.253 & 2.250 \\
Equal pairs & 32 & 2.595 & 2.562 & 2.586 \\
Equal pairs & 500 & 8.021 & 7.997 & 7.980 \\
Full overlap & 8 & 31.393 & 31.449 & 0.124 \\
Full overlap & 32 & limit & limit & 0.347 \\
Full overlap & 500 & limit & limit & 2.406 \\
Shifted 32 & 8 & 242.843 & 243.305 & 0.473 \\
Shifted 32 & 32 & limit & limit & 0.741 \\
Shifted 32 & 500 & limit & limit & 4.582 \\
Late defect & 8 & 33.184 & 33.148 & 32.865 \\
Late defect & 32 & limit & limit & limit \\
Late defect & 500 & limit & limit & limit \\
\bottomrule
\end{tabular}
\end{center}
Median milliseconds including construction; limits are censored. Full overlap enumerates all lengths in compressed progressions, not all occurrences.

\begin{center}\small
\begin{tabular}{@{}llrrr@{}}
\toprule
Family & $h$ & Before & A/A & Phases \\
\midrule
Partial & 8 & 0.581 & 0.586 & 0.585 \\
Partial & 32 & 1.165 & 1.200 & 1.193 \\
Partial & 500 & 7.128 & 7.086 & 7.116 \\
Equal-pair cyclic & 8 & 0.495 & 0.523 & 0.501 \\
Equal-pair cyclic & 32 & 0.832 & 0.846 & 0.838 \\
Equal-pair cyclic & 500 & 6.776 & 6.792 & 6.851 \\
Quotient & 8 & 0.348 & 0.351 & 0.340 \\
Quotient & 32 & 0.966 & 0.965 & 0.980 \\
Quotient & 500 & 13.381 & 13.702 & 13.413 \\
\bottomrule
\end{tabular}
\end{center}
Median milliseconds including construction; limits are censored. Local operations do not produce knot verdicts.


For the full-overlap family at $h=8$, median time falls from 31.393
milliseconds (31.449 A/A) to 0.124; charged work falls from 37641 to
328 and arena nodes from 108 to 11. At $h=32$ and $h=500$, both archived
controls exhaust two million work units, while phase certification completes
in 0.347 and 2.406 milliseconds respectively. The $h=500$ run uses 503
nodes, 4264 work units and no occurrence-table cells, returning one
progression representing $2^{500}$ overlap lengths.

For the shifted 32-letter block, at $h=8$ the medians are 242.843
milliseconds before, 243.305 A/A and 0.473 after. At $h=500$, only the
new stage completes: 4.582 milliseconds, 1094 nodes and 10110 charged
work units, again with one progression and no occurrence tables. This
confirms that the method handles a nonzero phase and is not restricted
to identical binary-power roots or a two-letter alphabet.

The late-defect control records the cost of an unsuccessful proposal.
At $h=8$, medians are 33.184, 33.148 and 32.865 milliseconds; the exact
work increase is 301 units, from 40113 to 40414. Nodes and cumulative
table cells remain 121 and 541 respectively. At $h=32$ and $h=500$,
all arms still exhaust the allowance. For $h=500$ the recorded times to
exhaustion are 2.717, 2.759 and 2.803 seconds. These are censored searches,
not completed-query timing ratios. The failed certificate correctly leaves
this more general obstacle unresolved.

All full knot queries produce identical certificate digests, and none
attempts the new phase stage. The small-fourteen sums of medians are
300.269 milliseconds before, 309.512 A/A and 299.232 after; Gordian
medians are 2.113, 2.229 and 2.141 seconds. Every local relator operation
also completes with the same trace across arms. Successful $h=500$
quotient-search medians are 13.381, 13.702 and 13.413 milliseconds.
These controls provide no evidence of a knot-corpus speedup, and their
A/A variation limits interpretation of small timing differences.

The archived baseline is the window-localized matcher at \code{ca4681932};
A/A repeats it. Every arm runs inside the same current arena, donor helpers,
recognizer and independent certificate replayer. Seed 3013 shuffles the
three arms in one excluded warmup and five measured rounds per case.
The audit includes 225 complete knot queries, 270 string queries and 135
local relator operations. String and relator timings include grammar
construction. Whole-query timings include fresh PD construction and replay,
with certificate digest comparison outside the clock. Local operations
have the unchanged two-million-work and 100000-node/table-cell allowances;
whole recognition has 18 seconds, including a 15-second, twenty-million-work
group stage. Limited runs are censored, not negative answers or completed
runtime denominators.

The late-defect control uses $X=Y=(ab)^Nc$. Its sampled prefix suggests
period two, but the final $c$ invalidates every phase mask. Its only overlap
is the full word; the complete fallback must establish this. This is a
correctness and speculation-cost control, not an expected shortcut success.

Raw samples are in \path{results/lcs_periodic_20261008.json}, and logs are
under \path{synthesis/data/lcs-periodic-*}. The incoming archives and their
relevance review are unchanged from Section~\ref{sec:lcs-windows}; the
suggested \code{openai/math} head remains \code{fd4aeeb2ee}. The phase
certificate is the elementary evaluation of a fixed periodic language on
a straight-line grammar, not a claim of a new general recognition theorem.
